\chapter{Lezione 6 -- Goldreich-Levin e cifratura computazionale}

\emph{Segue gli appunti del prof. Le parti aggiunte per capire meglio sono segnate ``non negli appunti del prof''. Usa i concetti della Lezione 4 (OWF, trascurabile) e della Lezione 5 (PRG, hard-core bit, riduzioni, ≈c).}

\section{Fine lezione 5: perché f\_BAD è comunque una OWF}

\begin{inbreve}

Nella lezione 5 abbiamo costruito \(f_{\text{BAD}}(x) = f(x) \,\|\, h(x)\), che scrive in chiaro il bit \(h(x)\): quindi \(h\) non è hard-core per \(f_{\text{BAD}}\). Qui il prof dimostra che \(f_{\text{BAD}}\) resta comunque \textbf{one-way}, con una riduzione: chi inverte \(f_{\text{BAD}}\) sa invertire anche \(f\).

\end{inbreve}

\textbf{Il problema}: \(h(\cdot)\) chiaramente \textbf{non} è hard-core per \(f_{\text{BAD}}\), perché è scritto nell'output. Perché \(f_{\text{BAD}}\) è ancora una OWF?

\textbf{Supponiamo di no}: esiste un PPT \(A\) che vince \(\text{GAME}^{\text{owf}}_{f_{\text{BAD}}}(\lambda)\) con probabilità \(\ge 1/\text{poly}(\lambda)\), cioè non trascurabile. Costruiamo un PPT \(B\) che \textbf{rompe} \(f\).

\begin{verbatim}
A                              B                              C
                                                              x ← U_n
                               ←────────── y = f(x) ──────────
                               b̃ ← U_1   (bit tirato a caso)
  ←──────── y ‖ b̃ ─────────
  ────────── x' ───────────→
                               ────────────── x' ────────────→
\end{verbatim}

\begin{enumerate}
\def\labelenumi{\arabic{enumi}.}
\tightlist
\item
  \(B\) riceve dal challenger \(y = f(x)\), con \(x\) casuale.
\item
  Ad \(A\) serve una stringa della forma \(f(x) \,\|\, h(x)\), ma \(B\) non conosce \(x\) e quindi non sa calcolare \(h(x)\). Allora \textbf{tira a indovinare}: sceglie un bit casuale \(\tilde{b}\) e dà ad \(A\) la stringa \(y \,\|\, \tilde{b}\).
\item
  \(A\) risponde con un \(x'\), e \(B\) lo inoltra al challenger.
\end{enumerate}

\textbf{Perché funziona} \emph{(spiegazione non negli appunti del prof)}:

\begin{itemize}
\tightlist
\item
  Con probabilità \(1/2\) il bit indovinato è giusto: \(\tilde{b} = h(x)\). In quel caso \(A\) riceve esattamente \(f_{\text{BAD}}(x)\), cioè la stessa cosa che riceverebbe nel suo gioco, quindi vince con la sua solita probabilità.
\item
  Se \(A\) vince, ha trovato un \(x'\) con \(f(x') \,\|\, h(x') = y \,\|\, \tilde{b}\). Guardando solo i primi bit: \(f(x') = y\). Quindi \(x'\) è una controimmagine di \(y\) e anche \(B\) vince.
\end{itemize}

\[
\Pr[B \text{ vince}] \ \ge\ \frac{1}{2} \cdot \Pr[A \text{ vince}] \ \ge\ \frac{1}{2} \cdot \frac{1}{\text{poly}(\lambda)} = \frac{1}{\text{poly}(\lambda)}
\]

Quindi \(B\) inverte \(f\) con probabilità non trascurabile, ma \(f\) è una OWF: \textbf{contraddizione}. ∎

\begin{nota}

\textbf{Il trucco da ricordare} \emph{(non negli appunti del prof)}: quando la riduzione non sa calcolare un pezzo dell'input dell'avversario, può \textbf{tirare a indovinare}. Se indovina con probabilità \(1/2\) (o \(1/\text{poly}\)), perde solo un fattore polinomiale, e poly · (1/poly) resta 1/poly.

\end{nota}

\section{Teorema di Goldreich-Levin}

\begin{inbreve}

Abbiamo visto che nessun \(h\) fissato è hard-core per \textbf{tutte} le OWF. La cosa migliore che si può avere è il contrario: \textbf{per ogni} OWF \(f\) esiste un hard-core bit. Il teorema lo costruisce: si aggiunge all'input una stringa casuale \(r\), la si pubblica, e l'hard-core bit è il \textbf{prodotto scalare modulo 2} tra \(x\) e \(r\).

\end{inbreve}

\emph{(Il nome ``Goldreich-Levin'' non è negli appunti del prof, che lo chiama solo THM.)}

\textbf{La cosa migliore possibile}: non ``\(\exists h\ \forall f\)'' (impossibile, per \(f_{\text{BAD}}\)), ma ``\(\forall f\ \exists h\)''.

\subsection{Prodotto scalare modulo 2}

\emph{(Definizione non negli appunti del prof.)}

Per due stringhe di bit \(x = (x_1, \dots, x_n)\) e \(r = (r_1, \dots, r_n)\):

\[
\langle x, r \rangle = \sum_{i=1}^{n} x_i r_i \bmod 2
\]

Si moltiplicano i bit a coppie (cioè si fa l'AND) e si guarda se il numero di 1 ottenuti è pari (risultato 0) o dispari (risultato 1). In pratica: è lo \textbf{XOR dei bit di} \(x\) \textbf{nelle posizioni in cui} \(r\) \textbf{vale 1}.

\begin{esempio}

\textbf{Esempio}: \(x = 1011\), \(r = 0110\).

\begin{itemize}
\tightlist
\item
  posizione 1: \(1 \cdot 0 = 0\)
\item
  posizione 2: \(0 \cdot 1 = 0\)
\item
  posizione 3: \(1 \cdot 1 = 1\)
\item
  posizione 4: \(1 \cdot 0 = 0\)
\end{itemize}

Somma \(= 1\), quindi \(\langle x, r \rangle = 1\). Detto in altro modo: \(r\) ``seleziona'' le posizioni 2 e 3 di \(x\), cioè i bit 0 e 1, e \(0 \oplus 1 = 1\).

\end{esempio}

\subsection{Enunciato (del prof)}

Sia \(f : \{0,1\}^n \to \{0,1\}^n\) una OWF. Consideriamo

\[
g : \{0,1\}^{2n} \to \{0,1\}^{2n} \qquad g(x, r) = f(x) \,\|\, r \qquad \forall\, x, r \in \{0,1\}^n
\]

Allora \(g\) è una OWF e

\[
h(x, r) = \langle x, r \rangle = \sum_{i=1}^{n} x_i r_i \bmod 2
\]

è \textbf{hard-core} per \(g\).

\begin{nota}

\textbf{Cosa sono le cose}

\begin{itemize}
\tightlist
\item
  \(x\) = l'input segreto della OWF \(f\); \(r\) = una stringa casuale in più, della stessa lunghezza.
\item
  \(g(x, r) = f(x) \,\|\, r\): calcola \(f(x)\) e poi scrive \(r\) \textbf{in chiaro}. Quindi chi vede l'output di \(g\) conosce \(r\) ma non \(x\).
\item
  \(h(x, r) = \langle x, r \rangle\): lo XOR di un \textbf{sottoinsieme casuale} dei bit di \(x\), quello scelto da \(r\).
\item
  In parole: anche conoscendo \(f(x)\) e \(r\), il bit \(\langle x, r \rangle\) sembra completamente casuale.
\end{itemize}

\end{nota}

\textbf{Perché} \(g\) \textbf{è una OWF} \emph{(non negli appunti del prof)}: invertire \(g\) vuol dire, dato \(f(x) \,\|\, r\), trovare \((x', r')\) con \(f(x') \,\|\, r' = f(x) \,\|\, r\). In particolare \(f(x') = f(x)\): chi inverte \(g\) inverte anche \(f\). È la stessa riduzione di \(f_{\text{BAD}}\), ma senza nemmeno dover indovinare (\(r\) è già dato).

\subsection{Intuizione (del prof)}

Supponiamo che esista un avversario \textbf{bravissimo} che, dato \(g(x, r)\), calcola \textbf{sempre} \(h(x, r)\), per ogni \(r\). Allora si ricava facilmente tutto \(x\), scegliendo come \(r\) le stringhe con un solo 1 (\(e_i\) = tutti 0 tranne un 1 in posizione \(i\)):

\[
h(x, e_1) = \langle x, 10\cdots0 \rangle = x_1 \qquad \cdots \qquad h(x, e_n) = \langle x, 0\cdots01 \rangle = x_n
\]

Con \(r = e_i\) il prodotto scalare ``seleziona'' solo il bit \(i\) di \(x\). Facendo \(n\) domande si ricostruisce \(x\) bit per bit, cioè si inverte \(f\): impossibile, perché \(f\) è one-way.

\begin{esempio}

\textbf{Esempio} \emph{(non negli appunti del prof)}: \(x = 101\). \(\langle x, 100 \rangle = 1\), \(\langle x, 010 \rangle = 0\), \(\langle x, 001 \rangle = 1\): si legge \(x = 101\).

\end{esempio}

\textbf{Perché la dimostrazione vera è più difficile}: un avversario reale indovina \(h(x, r)\) solo con probabilità \(\frac{1}{2} + \frac{1}{\text{poly}}\), e solo \textbf{in media su} \(r\) \textbf{casuale}. Sulle stringhe speciali \(e_i\) potrebbe sbagliare sempre. La dimostrazione completa deve ``amplificare'' questo piccolo vantaggio, e il prof non la fa.

\section{Usare i PRG per battere Shannon}

\begin{inbreve}

Shannon dice che con la segretezza perfetta la chiave deve essere lunga quanto il messaggio. Con un PRG si aggira il limite: si usa una chiave \textbf{corta}, la si allunga con il PRG e si usa il risultato come un One-Time Pad.

\end{inbreve}

\subsection{Lo schema (del prof)}

Con un PRG \(G : \{0,1\}^n \to \{0,1\}^{2n}\), chiave \(k \in \{0,1\}^n\) e messaggi \(m, c \in \{0,1\}^{2n}\):

\[
\text{Enc}(k, m) = G(k) \oplus m \qquad \text{Dec}(k, c) = G(k) \oplus c
\]

È un OTP in cui la chiave \(G(k)\) non è davvero casuale, ma \textbf{pseudocasuale}. La correttezza è la stessa dell'OTP: \(G(k) \oplus (G(k) \oplus m) = m\).

\textbf{La chiave è corta}: \(n\) bit per un messaggio di \(2n\) bit, quindi \textbf{più corta del messaggio}! Shannon non viene contraddetto, perché lo schema non è perfettamente segreto: è sicuro solo contro avversari efficienti.

\begin{esempio}

\textbf{Esempio} \emph{(valori di} \(G\) \emph{inventati, non negli appunti del prof)}: \(n = 3\), \(k = 101\), e supponiamo \(G(101) = 110010\). Per cifrare \(m = 011011\): \(c = 110010 \oplus 011011 = 101001\). Bob, che conosce \(k\), ricalcola \(G(101) = 110010\) e fa \(101001 \oplus 110010 = 011011 = m\). Si è usata una chiave di 3 bit per un messaggio di 6.

\end{esempio}

\subsection{È ancora one-time}

Come l'OTP, questo schema si può usare \textbf{una sola volta}. Dati \(c_1 = G(k) \oplus m_1\) e \(c_2 = G(k) \oplus m_2\):

\[
c_1 \oplus c_2 = m_1 \oplus m_2
\]

perché \(G(k) \oplus G(k) = 0\): la chiave sparisce, esattamente come nella Lezione 1.

\section{Sicurezza computazionale per la cifratura}

\begin{inbreve}

La segretezza perfetta chiede che i cifrati di due messaggi qualunque siano \textbf{identici} come distribuzione: nessuno, nemmeno con potenza infinita, li distingue. La versione computazionale chiede solo che siano \textbf{indistinguibili per avversari efficienti} (\(\approx_c\)).

\end{inbreve}

\textbf{Ripasso (segretezza perfetta)}: per ogni \(m_0, m_1\),

\[
\text{Enc}(K, m_0) \equiv \text{Enc}(K, m_1)
\]

quindi nemmeno un avversario \textbf{illimitato} riesce a distinguerli.

\subsection{Definizione: sicurezza computazionale one-time (del prof)}

Uno SKE \(\Pi = (\text{Enc}, \text{Dec})\) è \textbf{one-time computazionalmente sicuro} se

\[
\text{GAME}^{\text{one-time}}_{\Pi, A}(\lambda, 0) \ \approx_c\ \text{GAME}^{\text{one-time}}_{\Pi, A}(\lambda, 1)
\]

dove \(\text{GAME}^{\text{one-time}}_{\Pi, A}(\lambda, b)\), con \(b \in \{0, 1\}\), è:

\begin{verbatim}
A                                    C
  ────────── m0, m1 ∈ M ──────────→
                                     k ← K
                                     c = Enc(k, m_b)
  ←─────────────── c ──────────────
  ─────────────── b' ─────────────→
\end{verbatim}

\begin{enumerate}
\def\labelenumi{\arabic{enumi}.}
\tightlist
\item
  \(A\) sceglie \textbf{due messaggi} \(m_0, m_1\) e li manda al challenger.
\item
  Il challenger sceglie una chiave \(k\) a caso e cifra \(m_b\): \(c = \text{Enc}(k, m_b)\).
\item
  \(A\) vede \(c\) e risponde \(b'\), il suo tentativo di indovinare quale dei due è stato cifrato.
\end{enumerate}

In parole: \textbf{nessun avversario PPT} si comporta in modo (non trascurabilmente) diverso a seconda che sia stato cifrato \(m_0\) o \(m_1\).

\begin{nota}

\textbf{Differenze con il gioco della Lezione 2} \emph{(non negli appunti del prof)}:

\begin{itemize}
\tightlist
\item
  lì l'avversario poteva essere \textbf{illimitato} e la differenza doveva essere \textbf{esattamente 0}; qui l'avversario è \textbf{PPT} e la differenza deve essere \textbf{trascurabile};
\item
  qui è l'avversario a \textbf{scegliere} \(m_0\) e \(m_1\): la definizione deve reggere anche per i messaggi più comodi per lui.
\end{itemize}

\end{nota}

\subsection{Sanity checks (del prof)}

Una buona definizione deve implicare le cose che ci aspettiamo da un cifrario. Questa le implica:

\begin{itemize}
\tightlist
\item
  \textbf{È difficile calcolare la chiave} \(k\).

  \begin{itemize}
  \tightlist
  \item
    \textbf{Attenzione}: non basta! \(\text{Enc}(k, m) = m\) (non cifrare proprio) soddisfa anche lui la proprietà che \(k\) è difficile da calcolare, perché il cifrato non dice niente della chiave. Ma è ovviamente insicuro.
  \end{itemize}
\item
  \textbf{È difficile calcolare il messaggio} \(m\).
\item
  \textbf{Ma anche} il primo bit di \(m\), il secondo bit, \ldots{} : nessuna informazione parziale sul messaggio.
\end{itemize}

\emph{Perché la definizione implica tutto questo (non negli appunti del prof)}: se un avversario sapesse calcolare \(k\) (o \(m\), o il primo bit di \(m\)), vincerebbe il gioco. Basta scegliere \(m_0\) e \(m_1\) che differiscono proprio in quello che sa calcolare (es. \(m_0\) che inizia con 0 e \(m_1\) che inizia con 1), calcolarlo dal cifrato e rispondere di conseguenza.

\section{Teorema: il PRG dà cifratura one-time sicura}

\begin{inbreve}

Se \(G\) è un PRG, lo schema \(\text{Enc}(k, m) = G(k) \oplus m\) è one-time computazionalmente sicuro. La dimostrazione usa il \textbf{game hopping}: si passa dal gioco con \(b = 0\) a quello con \(b = 1\) attraverso giochi intermedi, mostrando che ogni salto è invisibile.

\end{inbreve}

\subsection{Enunciato (del prof)}

Assumendo che \(G\) sia un PRG, lo SKE definito sopra è \textbf{one-time computazionalmente sicuro}.

\subsection{Game hopping}

\textbf{Game hopping} (tecnica di dimostrazione molto comune): si parte dai due giochi da confrontare,

\[
G_0 \equiv \text{GAME}^{\text{one-time}}_{\Pi, A}(\lambda, 0) \qquad G_1 \equiv \text{GAME}^{\text{one-time}}_{\Pi, A}(\lambda, 1)
\]

e bisogna mostrare \(G_0 \approx_c G_1\). Si introducono \textbf{due giochi in più}, \(H_0\) e \(H_1\):

\begin{verbatim}
                 G_b  /  H_b
A                                    C
  ────────── m0, m1 ──────────────→
                                     k ← {0,1}^n ;   z ← U_2n
                                     G_b:  c_b = G(k) ⊕ m_b
                                     H_b:  c_b = z ⊕ m_b
  ←─────────────── c ──────────────
  ─────────────── b' ─────────────→
\end{verbatim}

\begin{itemize}
\tightlist
\item
  In \(G_b\) si cifra con il PRG: \(c = G(k) \oplus m_b\) (lo schema vero).
\item
  In \(H_b\) si cifra con una stringa \textbf{davvero casuale} \(z\) di \(2n\) bit: \(c = z \oplus m_b\). È un vero OTP.
\end{itemize}

È la stessa idea dell'argomento ibrido della Lezione 5 \emph{(collegamento non negli appunti del prof)}: si sostituisce il pezzo pseudocasuale con uno davvero casuale.

\subsection{Il piano}

\begin{enumerate}
\def\labelenumi{\arabic{enumi}.}
\tightlist
\item
  \(H_0 \equiv H_1\): segue dalla \textbf{segretezza perfetta dell'OTP} (Lezione 1). Con \(z\) uniforme, \(z \oplus m_0\) e \(z \oplus m_1\) hanno la stessa distribuzione, quindi \(A\) si comporta in modo identico.
\item
  Resta da mostrare \(G_b \approx_c H_b\) (per ogni \(b\)). Allora:
\end{enumerate}

\[
G_0 \ \approx_c\ H_0 \ \equiv\ H_1 \ \approx_c\ G_1
\]

\emph{Perché la catena basta (non negli appunti del prof)}: il vantaggio di \(A\) tra \(G_0\) e \(G_1\) è al più la somma dei vantaggi sui singoli passi (disuguaglianza triangolare): \(\text{negl} + 0 + \text{negl} = \text{negl}\).

\subsection{La riduzione: G\_b ≈c H\_b}

\textbf{Supponiamo di no}: esiste un PPT \(A\) che distingue \(G_b\) da \(H_b\) (per qualche \(b \in \{0, 1\}\)) con probabilità \(\ge 1/\text{poly}\). Costruiamo \(B\) che rompe il PRG:

\begin{verbatim}
A                         B                              C
                                                         z = G(k), k ← U_n
                                                         oppure  z ← U_2n
                          ←──────────── z ─────────────
  ←────── "START" ──────
  ─────── m0, m1 ───────→
  ←──── c = m_b ⊕ z ─────
  ────────── b' ────────→
                          ──────────── b' ─────────────→
\end{verbatim}

\begin{enumerate}
\def\labelenumi{\arabic{enumi}.}
\tightlist
\item
  \(B\) riceve dal challenger del PRG una stringa \(z\) di \(2n\) bit: o \(z = G(k)\) con \(k\) casuale, o \(z\) davvero casuale.
\item
  \(B\) fa partire \(A\) (``START'') e riceve i suoi due messaggi \(m_0, m_1\).
\item
  \(B\) risponde con \(c = m_b \oplus z\) (usa \(z\) come se fosse la chiave dell'OTP).
\item
  Quando \(A\) risponde \(b'\), \(B\) inoltra \(b'\) al challenger.
\end{enumerate}

\textbf{Riduzione perfetta}: \(B\) simula esattamente uno dei due giochi.

\[
\Pr[B \Rightarrow 1 : z = G(k)] = \Pr[A \Rightarrow 1 : \text{in } G_b]
\]

\[
\Pr[B \Rightarrow 1 : z \leftarrow U_{2n}] = \Pr[A \Rightarrow 1 : \text{in } H_b]
\]

\begin{itemize}
\tightlist
\item
  Se \(z = G(k)\), allora \(c = G(k) \oplus m_b\): è esattamente \(G_b\).
\item
  Se \(z\) è casuale, allora \(c = z \oplus m_b\): è esattamente \(H_b\).
\end{itemize}

Quindi il vantaggio di \(B\) contro il PRG è uguale a quello di \(A\) tra \(G_b\) e \(H_b\), cioè \(\ge 1/\text{poly}\): \(B\) rompe \(G\). Ma \(G\) è un PRG: \textbf{contraddizione}. Inoltre \(B\) è PPT, perché fa solo uno XOR in più rispetto ad \(A\). ∎

\section{Riepilogo}

\begin{itemize}
\tightlist
\item
  \textbf{f\_BAD è una OWF}: la riduzione \(B\) non conosce \(h(x)\) e lo \textbf{indovina} con un bit casuale; \(\Pr[B \text{ vince}] \ge \frac{1}{2} \Pr[A \text{ vince}]\)
\item
  Nessun \(h\) va bene per tutte le \(f\), ma \textbf{per ogni} \(f\) \textbf{esiste} un \(h\) (\textbf{Goldreich-Levin}): \(g(x, r) = f(x) \,\|\, r\) è una OWF e \(h(x, r) = \langle x, r \rangle = \sum x_i r_i \bmod 2\) è hard-core per \(g\)
\item
  Intuizione: un avversario perfetto con \(r = e_i\) ricaverebbe \(x\) bit per bit; la dimostrazione vera è difficile perché l'avversario ha solo un vantaggio \(1/\text{poly}\) su \(r\) casuale
\item
  \textbf{Battere Shannon con un PRG}: \(\text{Enc}(k, m) = G(k) \oplus m\) con \(G : \{0,1\}^n \to \{0,1\}^{2n}\): chiave più corta del messaggio; ma è ancora \textbf{one-time} (\(c_1 \oplus c_2 = m_1 \oplus m_2\))
\item
  \textbf{Sicurezza one-time computazionale}: \(A\) sceglie \(m_0, m_1\), riceve \(\text{Enc}(k, m_b)\); \(\text{GAME}(\lambda, 0) \approx_c \text{GAME}(\lambda, 1)\). Implica che \(k\), \(m\) e ogni bit di \(m\) sono difficili da calcolare (``\(k\) difficile'' da solo non basta: vedi \(\text{Enc}(k, m) = m\))
\item
  \textbf{Teorema}: \(G\) PRG ⇒ lo schema è one-time computazionalmente sicuro. Prova con \textbf{game hopping}: \(G_0 \approx_c H_0 \equiv H_1 \approx_c G_1\), dove \(H_b\) usa un \(z\) davvero casuale (OTP) e \(G_b \approx_c H_b\) per riduzione perfetta al PRG
\end{itemize}
